% answers_short.tex —— 2012 年·广东卷·理科 · 简短答案速查
% 由 tools/make_short_answers.py 自 answers.tex 的 \gksolhead/\gkanswer 提取，
% 并与 2012/2012_数学真题_广东理_答案要点.md 逐题比对后生成，勿手改（改源 MD 后重跑）。

\begin{gknotice}
\textbf{使用说明}：本节只给最终答案便于快速核对，完整推导见解析版。答案为编写者独立推导，\textbf{非官方标准答案}。
\end{gknotice}

\gkpart{选择题}
\begin{center}\begin{tabular}{ccccccccc}
\toprule
题号 & 1 & 2 & 3 & 4 & 5 & 6 & 7 & 8 \\
\midrule
答案 & D & C & A & A & B & C & D & C \\
\bottomrule
\end{tabular}\end{center}

\gkpart{填空题}
\begin{itemize}[leftmargin=2.6em]
  \item[\textbf{9.}] $\left(-\infty,-\dfrac12\right]$。分 $x<-2$、$-2\leqslant x<0$、$x\geqslant0$ 三段讨论，仅前两段的部分满足。
  \item[\textbf{10.}] $20$。通项为 $\mathrm C_6^k x^{12-3k}$，令 $12-3k=3$ 得 $k=3$，系数 $\mathrm C_6^3=20$。
  \item[\textbf{11.}] $a_n=2n-1$。设公差 $d>0$，由 $1+2d=(1+d)^2-4$ 得 $d=2$。
  \item[\textbf{12.}] $y=2x+1$。$y'=3x^2-1$，在 $x=1$ 处斜率为 $2$。
  \item[\textbf{13.}] $8$。依次得 $s=2,4,8$，对应 $i=2,4,6$；$i=8$ 时退出循环。
  \item[\textbf{14.}] $(1,1)$。$C_1$ 为 $y^2=x$（$y\geqslant0$），$C_2$ 为 $x^2+y^2=2$，联立得 $x=1$。
  \item[\textbf{15.}] $\sqrt3$。$\angle AOC=2\angle ABC=60^\circ$，在直角三角形 $OAP$ 中 $PA=OA\tan60^\circ=\sqrt3$。
\end{itemize}

\gkpart{解答题}
\begin{itemize}[leftmargin=2.6em]
  \item[\textbf{16.}] (1) $\omega=\dfrac15$。由 $\dfrac{2\pi}{\omega}=10\pi$ 得；(2) $-\dfrac{13}{85}$。由条件得 $\sin\alpha=\dfrac35$，$\cos\alpha=\dfrac45$，$\cos\beta=\dfrac8{17}$，$\sin\beta=\dfrac{15}{17}$，再用两角和的余弦公式
  \item[\textbf{17.}] (1) $x=0.018$。各组频率之和为 $1$，即 $10(0.082+x)=1$；(2) $\dfrac12$。不低于 $80$ 分者共 $12$ 人，其中不低于 $90$ 分者 $3$ 人，故 $E\xi=2\cdot\dfrac3{12}=\dfrac12$
  \item[\textbf{18.}] (1) 见解析版；(2) $3$。由 $PC\perp BD$ 定出 $AB=2$，两半平面法向量可取 $(1,0,2)$、$(-1,1,0)$，夹角余弦 $\dfrac1{\sqrt{10}}$，正切为 $3$
  \item[\textbf{19.}] (1) $a_1=1$。由 $n=1,2$ 得 $a_2=2a_1+3$、$a_3=2a_1+2a_2+7$，代入等差条件即得；(2) $a_n=3^n-2^n$。作差得 $a_{n+1}=3a_n+2^n$，化归等比数列求解；(3) 见解析版
  \item[\textbf{20.}] (1) $\dfrac{x^2}3+y^2=1$。$e^2=\dfrac23$ 给 $b^2=\dfrac{a^2}3$，距离平方的最大值为 $a^2+6=9$；(2) 存在，$M$ 为 $\left(\pm\dfrac{\sqrt6}2,\pm\dfrac{\sqrt2}2\right)$（四个点），$\triangle OAB$ 面积的最大值为 $\dfrac12$。面积取最大值时圆心到直线距离为 $\dfrac1{\sqrt2}$，即 $m^2+n^2=2$，再与椭圆方程联立
  \item[\textbf{21.}] (1) 记 $\Delta=9a^2-30a+9=3(3a-1)(a-3)$，$x_{1,2}=\dfrac{3(1+a)\mp\sqrt\Delta}4$（$x_1<x_2$）。当 $a<0$ 时 $D=(x_2,+\infty)$；当 $a=0$ 时 $D=\left(\dfrac32,+\infty\right)$；当 $0<a<\dfrac13$ 时 $D=(0,x_1)\cup(x_2,+\infty)$；当 $a=\dfrac13$ 时 $D=(0,1)\cup(1,+\infty)$；当 $\dfrac13<a<1$ 时 $D=(0,+\infty)$；(2) $a\leqslant0$ 时无极值点；$0<a<\dfrac13$ 时极大值点为 $x=a$；$a=\dfrac13$ 时极大值点为 $x=\dfrac13$；$\dfrac13<a<1$ 时极大值点为 $x=a$、极小值点为 $x=1$。依据 $f'(x)=6(x-1)(x-a)$ 及 $a$、$1$ 是否落在 $D$ 内判定
\end{itemize}
